% Options for packages loaded elsewhere
% Options for packages loaded elsewhere
\PassOptionsToPackage{unicode}{hyperref}
\PassOptionsToPackage{hyphens}{url}
\PassOptionsToPackage{dvipsnames,svgnames,x11names}{xcolor}
%
\documentclass[
  letterpaper,
]{article}
\usepackage{xcolor}
\usepackage{amsmath,amssymb}
\setcounter{secnumdepth}{5}
\usepackage{iftex}
\ifPDFTeX
  \usepackage[T1]{fontenc}
  \usepackage[utf8]{inputenc}
  \usepackage{textcomp} % provide euro and other symbols
\else % if luatex or xetex
  \usepackage{unicode-math} % this also loads fontspec
  \defaultfontfeatures{Scale=MatchLowercase}
  \defaultfontfeatures[\rmfamily]{Ligatures=TeX,Scale=1}
\fi
\usepackage{lmodern}
\ifPDFTeX\else
  % xetex/luatex font selection
\fi
% Use upquote if available, for straight quotes in verbatim environments
\IfFileExists{upquote.sty}{\usepackage{upquote}}{}
\IfFileExists{microtype.sty}{% use microtype if available
  \usepackage[]{microtype}
  \UseMicrotypeSet[protrusion]{basicmath} % disable protrusion for tt fonts
}{}
\makeatletter
\@ifundefined{KOMAClassName}{% if non-KOMA class
  \IfFileExists{parskip.sty}{%
    \usepackage{parskip}
  }{% else
    \setlength{\parindent}{0pt}
    \setlength{\parskip}{6pt plus 2pt minus 1pt}}
}{% if KOMA class
  \KOMAoptions{parskip=half}}
\makeatother
% Make \paragraph and \subparagraph free-standing
\makeatletter
\ifx\paragraph\undefined\else
  \let\oldparagraph\paragraph
  \renewcommand{\paragraph}{
    \@ifstar
      \xxxParagraphStar
      \xxxParagraphNoStar
  }
  \newcommand{\xxxParagraphStar}[1]{\oldparagraph*{#1}\mbox{}}
  \newcommand{\xxxParagraphNoStar}[1]{\oldparagraph{#1}\mbox{}}
\fi
\ifx\subparagraph\undefined\else
  \let\oldsubparagraph\subparagraph
  \renewcommand{\subparagraph}{
    \@ifstar
      \xxxSubParagraphStar
      \xxxSubParagraphNoStar
  }
  \newcommand{\xxxSubParagraphStar}[1]{\oldsubparagraph*{#1}\mbox{}}
  \newcommand{\xxxSubParagraphNoStar}[1]{\oldsubparagraph{#1}\mbox{}}
\fi
\makeatother


\usepackage{longtable,booktabs,array}
\usepackage{calc} % for calculating minipage widths
% Correct order of tables after \paragraph or \subparagraph
\usepackage{etoolbox}
\makeatletter
\patchcmd\longtable{\par}{\if@noskipsec\mbox{}\fi\par}{}{}
\makeatother
% Allow footnotes in longtable head/foot
\IfFileExists{footnotehyper.sty}{\usepackage{footnotehyper}}{\usepackage{footnote}}
\makesavenoteenv{longtable}
\usepackage{graphicx}
\makeatletter
\newsavebox\pandoc@box
\newcommand*\pandocbounded[1]{% scales image to fit in text height/width
  \sbox\pandoc@box{#1}%
  \Gscale@div\@tempa{\textheight}{\dimexpr\ht\pandoc@box+\dp\pandoc@box\relax}%
  \Gscale@div\@tempb{\linewidth}{\wd\pandoc@box}%
  \ifdim\@tempb\p@<\@tempa\p@\let\@tempa\@tempb\fi% select the smaller of both
  \ifdim\@tempa\p@<\p@\scalebox{\@tempa}{\usebox\pandoc@box}%
  \else\usebox{\pandoc@box}%
  \fi%
}
% Set default figure placement to htbp
\def\fps@figure{htbp}
\makeatother





\setlength{\emergencystretch}{3em} % prevent overfull lines

\providecommand{\tightlist}{%
  \setlength{\itemsep}{0pt}\setlength{\parskip}{0pt}}



 


\usepackage{orcidlink}  
\usepackage[nonatbib, preprint]{neurips_2023}
\usepackage{amsmath}
\usepackage{amsthm}
\usepackage[utf8]{inputenc} % allow utf-8 input
\usepackage[T1]{fontenc}    % use 8-bit T1 fonts
\usepackage{hyperref}       % hyperlinks
\usepackage{url}            % simple URL typesetting
\usepackage{booktabs}       % professional-quality tables
\usepackage{amsfonts}       % blackboard math symbols
\usepackage{nicefrac}       % compact symbols for 1/2, etc.
\usepackage{microtype}      % microtypography
\usepackage{xcolor}         % colors
\usepackage[numbers]{natbib}
\usepackage{comment}
\usepackage{graphicx} 
\bibliographystyle{abbrvnat}
\makeatletter
\@ifpackageloaded{caption}{}{\usepackage{caption}}
\AtBeginDocument{%
\ifdefined\contentsname
  \renewcommand*\contentsname{Table of contents}
\else
  \newcommand\contentsname{Table of contents}
\fi
\ifdefined\listfigurename
  \renewcommand*\listfigurename{List of Figures}
\else
  \newcommand\listfigurename{List of Figures}
\fi
\ifdefined\listtablename
  \renewcommand*\listtablename{List of Tables}
\else
  \newcommand\listtablename{List of Tables}
\fi
\ifdefined\figurename
  \renewcommand*\figurename{Figure}
\else
  \newcommand\figurename{Figure}
\fi
\ifdefined\tablename
  \renewcommand*\tablename{Table}
\else
  \newcommand\tablename{Table}
\fi
}
\@ifpackageloaded{float}{}{\usepackage{float}}
\floatstyle{ruled}
\@ifundefined{c@chapter}{\newfloat{codelisting}{h}{lop}}{\newfloat{codelisting}{h}{lop}[chapter]}
\floatname{codelisting}{Listing}
\newcommand*\listoflistings{\listof{codelisting}{List of Listings}}
\makeatother
\makeatletter
\makeatother
\makeatletter
\@ifpackageloaded{caption}{}{\usepackage{caption}}
\@ifpackageloaded{subcaption}{}{\usepackage{subcaption}}
\makeatother
\usepackage{bookmark}
\IfFileExists{xurl.sty}{\usepackage{xurl}}{} % add URL line breaks if available
\urlstyle{same}
\hypersetup{
  pdftitle={Measuring Alpha Decay: Analysis of Order Splitting via Trade Markouts},
  pdfauthor={John Joe; John Joe; John Joe},
  pdfkeywords={template, quarto, R, Python, reproductibility},
  colorlinks=true,
  linkcolor={blue},
  filecolor={Maroon},
  citecolor={Blue},
  urlcolor={Blue},
  pdfcreator={LaTeX via pandoc}}



\title{Measuring Alpha Decay: Analysis of Order Splitting via Trade
Markouts}
\author{
\textbf{John Joe}%
~\orcidlink{0000-0000-0000-0000}\\%
%
    Electrical and Computer Engineering\\%
    Texas A\&M University\\%
    College Station, TX 77843, USA\\%
%
\texttt{computo@sfds.asso.fr}\and 
\textbf{John Joe}%
\\%
%
\and 
\textbf{John Joe}%
\\%
%
}
\date{2026-04-07}
\begin{document}
\maketitle
\begin{abstract}
A fundamental hurdle in trading is the balance between market impact and
alpha decay. Submitting a large order leads to significant price impact
as it takes the liquidity from the Limit Order Book (LOB). On the other
hand, splitting a large order into smaller chunks over seconds or
minutes exposes unexecuted orders to alpha decay and price drifts. This
hurdle is especially prevalent in crowded markets where thousands of
algorithms are trading on similar signals.

The goal of the project will be to measure the average alpha lost when
forced to split orders by comparing different execution strategies. In
the execution of this goal, will use post-trade markouts to evaluate the
price impact and the drift over short periods of time as we follow the
signal. Additionally, we will record the average cost of entering the
position, the aim of which is to minimize the cost.
\end{abstract}


\section{Goals for the Milestone}\label{goals-for-the-milestone}

The goals we've set for this milestone report are as follows.

\begin{enumerate}
\def\labelenumi{\arabic{enumi}.}
\tightlist
\item
  A data ingestion pipeline, which will parse L2 order book data from
  Databento for use in our \texttt{ExecutionSimulator}.
\item
  An \texttt{ExecutionSimulator} which can simulate the effects and
  performance of multiple execution strategies, as well as the effects
  of latency and LOB replenishment on the fill strategies.
\item
  Markout plots, which help inform the execution strategies. We use
  markouts to explore the effectiveness of order book imbalance (OBI) as
  a predictor of short term price movements that may inform an execution
  strategy.
\end{enumerate}

\section{Developments}\label{developments}

\subsection{Sourcing Data}\label{sourcing-data}

We are using the \texttt{MBP-10} schema from Databento, which provides
the top 10 levels of the Limit Order Book (LOB) for a selection of
stocks and date range. The dataset includes all order book updates
(trades, adds, cancels, etc.) which allows us to reconstruct the market
dynamics at a granular level. The data is sourced from
\href{https://databento.com/docs/schemas-and-data-formats/mbp-10?historical=python&live=python&reference=python}{Databento}.

In the project milestone, we utilize L2 data for \texttt{AMZN} from
\texttt{"2026-03-25":"2026-03-27"}. We are intentionally using a small
time range and one ticker in order to minimize runtime and memory
requirements while we develop the methodology we need. Through the end
of the project, we will be adding multiple assets spanning the
large/small market cap and tick-size range to explore how different
assets require different execution strategies.

Mechanically, we have written a simple \texttt{MBP10DataAdapter} class
that uses a hierarchical cacheing structure. This allows us to fetch
data from Databento only when absolutely necessary, which minimizes cost
and tiem spent downloading data. The adapter caches on
\texttt{{[}tickers,\ start,\ end,\ dataset,\ schema{]}}. Here, our
schema is always \texttt{MBP-10} and we use the \texttt{XNAS.ITCH}
Dataset which gives us access to only the Nasdaq exchange.

\section{Findings}\label{findings}

\section{Goals for the Future}\label{goals-for-the-future}




\end{document}
